% Section content template for individual Educative lessons
% This file contains the actual content for a specific section
% Generated from Educative API section content

% Section: Code for Stack Overflow
% Section ID: 5913671215874048
% Section Slug: code-for-stack-overflow
% Generated: 2025-12-12T22:50:28.857135

We've gone over the different aspects of Stack Overflow and observed the attributes attached to the problem using various UML diagrams. Let us now explore the more practical side of things, where we will work on implementing the Stack Overflow network using multiple languages. This is usually the last step in an object-oriented design interview process.
\\
We have chosen the following languages to write the skeleton code of the different classes present in Stack Overflow:

\begin{itemize}
\item
\protect\phantomsection\label{FM09jbN_gaYyJplkwgTn3}
Java
\item
\protect\phantomsection\label{9KiWz9epPYOCf0dhSwh5C}
C\#
\item
\protect\phantomsection\label{7NQBsSvUXqoSYQLlf7sJX}
Python
\item
\protect\phantomsection\label{rOSnN5aLDXv5BltyoRmQW}
C++
\item
\protect\phantomsection\label{5tqJBV2aVwS62Og1eMLN1}
JavaScript
\end{itemize}
\\
If you want to skip directly to the implementation and see the complete workflow, simply scroll down or click~\hyperref[Executable-code-Designing-Stack-Overflow]{Executable Code: Designing Stack Overflow}.

\subsection{Stack Overflow classes}\label{xJGQXkMtzc9iBGZqUgigK}
\\
This section will provide the skeleton code of the classes designed in the class diagram lesson.
\begin{quote}
\textbf{Note:} For simplicity, we are not defining getter and setter functions. The reader can assume that all class attributes are private, accessed through their respective public getter methods, and modified only through their public method functions.
\end{quote}
\subsubsection{Constants}\label{YlMdA5Xd5mvd7O_iwtTJn}
\\
The following code defines the various enums and custom data types being used in the Stack Overflow design:
\begin{quote}
\textbf{Note:} JavaScript does not support enumerations, so we will use the \texttt{Object.freeze()} method as an alternative that freezes an object and prevents further modifications.
\end{quote}

\textbf{Constant definitions}

\begin{lstlisting}[language=Java]
enum QuestionStatus {
  ACTIVE, 
  CLOSED,
  FLAGGED,
  BOUNTIED
}

enum UserStatus {
  ACTIVE,
  BLOCKED
}
\end{lstlisting}

\subsubsection{User, admin, moderator, and guest}\label{a-GjKp6-VT6SMnzmDC3jh}
\\
The \texttt{User} class will be a parent class that represents a regular Stack Overflow user. A normal user can also be an \texttt{Admin} and a \texttt{Moderator}. Another actor is represented by the \texttt{Guest} class, which refers to a user who can only search and view questions and their answers. However, they need to register an account to ask or answer questions. The definition of these classes is provided below:

\textbf{The User, Admin, Moderator, and Guest classes}

\begin{lstlisting}[language=Java]
public class User {
  private int reputationPoints;
  private String username;
  private List<Badge> badges;
  private Question currentDraft;

  public Question create(String title, String body, List<Tag> tags);
  public boolean addAnswer(Question, question, Answer answer);
  public boolean createComment(Comment comment);
  public boolean createTag(Tag tag);
  public void flagQuestion(Question question);
  public void flagAnswer(Answer answer);
  public void upvote(int id);
  public void downvote(int id);
  public void voteToCloseQuestion(Question question);
  public void voteToDeleteQuestion(Question question);
  public void acceptAnswer(Answer answer);
  public boolean publishQuestion();
  public void addBounty(int value);
}

public class Admin extends User {
  public boolean blockUser(User user);
  public boolean unblockUser(User user);
  public void awardBadge(User user, Badge badge);
}

public class Moderator extends User {
  public void closeQuestion(Question question);
  public void reopenQuestion(Question question);
  public void deleteQuestion(Question question);
  public void restoreQuestion(Question question);
  public void deleteAnswer(Answer answer);
}

public class Guest {
  public void registerAccount();
  public List<Question> searchQuestions();
  public void viewQuestion(Question question);
}
\end{lstlisting}

\subsubsection{Question, answer, comment, and bounty}\label{3VyJ6YosyGRVL7Tb2VQ2W}
\\
Stack Overflow users can create and answer questions, upvote and downvote them, and add bounties and comments to questions. The definition of these classes is provided below:

\textbf{The Question, Answer, Comment, and Bounty classes}

\begin{lstlisting}[language=Java]
public class Question {
  private int id;
  private String title;
  private String content;
  private User createdBy;
  private QuestionStatus status;
  private int upvotes;
  private int downvotes;
  private int voteCount;
  private Date creationDate;
  private Date modificationDate;
  private Bounty bounty;

  private List<Tag> tags;
  private List<Comment> comments;
  private List<Answer> answers;
  private List<User> followers;

  public void addComment(Comment comment);
  public void addBounty(Bounty bounty);
  public void notify(String closedQuestion);
  public void createBounty(int value);
}

public class Comment {
  private int id;
  private String content;
  private int flagCount;
  private int upvotes;
  private Date creationDate;
  private User postedBy;
}

public class Answer {
  private int id;
  private String content;
  private int flagCount;
  private int voteCount;
  private int upvotes;
  private int downvotes;  
  private boolean isAccepted;
  private Date creationTime;
  private User postedBy;
  private List<Comment> comments;

  public void addComment(Comment comment);
}

public class Bounty {
  private int reputationPoints;
  private Date expiryDate;
  public boolean updateReputationPoints(int reputation);
}
\end{lstlisting}

\subsubsection{Badge, tag, and tag list}\label{julu7b82zRwsJjYIajaE4-1}
\\
Users can have badges that act as their reputation awards. Questions can have tags that describe the category that the question falls under. To keep a count of the tags being used, the \texttt{TagList} class is used. The definition of these classes is provided below:

\textbf{The Badge, Tag, and TagList classes}

\begin{lstlisting}[language=Java]
public class Badge {
  private String name;
  private String description;
}

public class Tag {
  private String name;
  private String description;
}

public class TagList {
  private HashMap<Tag, int> tagsCount;
  public void incrementTagCount(Tag tag);
  public void decrementTagCount(Tag tag);
}
\end{lstlisting}

\subsubsection{Notification}\label{0ZVCay7W6Bkab1KXP_cLz}
\\
The \texttt{Notification} class is responsible for sending users notifications about new messages, comments, posts, or friend requests via a phone number or email. Its definition is provided below:

\textbf{The Notification class}

\begin{lstlisting}[language=Java]
public class Notification {
  private int notificationId;
  private Date createdOn;
  private String content;

  public boolean sendNotification(User user);
}
\end{lstlisting}

\subsubsection{Search catalog and interface}\label{julu7b82zRwsJjYIajaE4-2}
\\
The \texttt{SearchCatalog} class contains information on existing questions and answers. It also implements the \texttt{Search} interface class to enable the search functionality based on the given criteria (tags, usernames, and searched keywords). The definition of these two classes is provided below:

\textbf{The Search interface and the SearchCatalog class}

\begin{lstlisting}[language=Java]
public interface Search {
    List<Question> searchByTags(String tagName);
    List<Question> searchByUsers(String username);
    List<Question> searchByWords(String word);
}

public class SearchCatalog implements Search {
    // Map tag names to lists of questions associated with those tags
    private HashMap<Tag, List<Question>> questionsUsingTags;
    // Map usernames to lists of questions created by those users
    private HashMap<User, List<Question>> questionsUsingUsers;
    // Map words to lists of questions containing those words
    private HashMap<String, List<Question>> questionsUsingWords;

    public List<Question> searchByTags(String tagName);
    public List<Question> searchByUsers(String username);
    public List<Question> searchByWords(String word);
    public void indexQuestion(Question question);
}
\end{lstlisting}

\subsection{Executable code: Designing Stack Overflow}\label{mRiWAVwGAjXzjjLIri0XA}
\\
Below is a fully self-contained, runnable program in Java, C\#, C++, Python, and JavaScript that demonstrates the core workflows of the Designing Stack Overflow system. The system's main driver code resides in~\texttt{Driver.java},~\texttt{Driver.cs},~\texttt{Driver.py},~\texttt{Driver.cpp}, and~\texttt{Driver.js}~for all respective languages. You can click the ``Run'' button to execute the code.

\subsubsection{What does this code show}\label{4VuIFXTFaihKUO3kDpZoP}

\begin{itemize}
\item
\protect\phantomsection\label{lZSH5-mzHIs8g6dW6Lcxy}
\textbf{System initialization:} Initializes a guest account and two users---Alice (\texttt{alice123}) and Bob (bob\_dev)---with assigned reputation points and active status. This sets up the environment for simulating Q\&A interactions on the platform.
\item
\protect\phantomsection\label{Bj-rjzM7H9hSbNYXd1CAU}
\textbf{Scenario 1 (Tag creation and question drafting):} Alice creates two tags, ``java'' and ``oop,'' and uses them to draft a question about the difference between interfaces and abstract classes in Java. This models how users contribute categorized content to the platform.
\item
\protect\phantomsection\label{QDaOr_ClfBb-eqXXd2527}
\textbf{Scenario 2 (Bounty addition and question publishing):} Alice adds a 100-point reputation bounty to encourage quality responses after drafting the question. She then publishes the question, simulating the transition from draft to public visibility.
\item
\protect\phantomsection\label{O7o2T6NduMuHyEo1DAlqH}
\textbf{Scenario 3 (Answer submission and acceptance):} Bob writes an answer to Alice's question, clearly distinguishing interfaces from abstract classes. Alice accepts Bob's answer, showcasing how accepted answers promote helpful contributions within the community.
\end{itemize}
\begin{quote}
\textbf{Hint: Try customizing the code!}\\
This code is designed for experimentation, not just reading. You can edit, extend, or modify any part of the example.
\end{quote}

\textbf{Executable Code: Designing Stack Overflow}

\begin{lstlisting}[language=Java]
import java.util.*;

public class Driver {
    public static void main(String[] args) {
        System.out.println("=== Welcome to Stack Overflow Simulation ===
");

        // ================================
        // ✅ System Initialization
        // ================================
        Guest guest = new Guest();
        guest.registerAccount();
        System.out.println("System initialized with a guest account.
");

        // ================================
        // 👤 Scenario 1: User Registration and Tag Creation
        // ================================

        // Creating two users: Alice and Bob
        User user1 = new User();
        user1.setReputationPoints(10);
        user1.setName("alice123");
        user1.setStatus(UserStatus.ACTIVE);

        User user2 = new User();
        user2.setReputationPoints(7);
        user2.setName("bob_dev");
        user2.setStatus(UserStatus.ACTIVE);

        System.out.println("=== Scenario 1: Two users registered. ===");
        System.out.println("- User1: Alice (alice123), Reputation: 10");
        System.out.println("- User2: Bob (bob_dev), Reputation: 7
");

        // Alice creates tags: java and oop
        Tag javaTag = new Tag();
        javaTag.setName("java");
        javaTag.setDescription("Java programming language");

        Tag oopTag = new Tag();
        oopTag.setName("oop");
        oopTag.setDescription("Object-oriented programming");

        List<Tag> tags = new ArrayList<>();
        tags.add(javaTag);
        tags.add(oopTag);

        System.out.println("Alice created two tags: [java, oop]
");

        // ================================
        // ❓ Scenario 2: Question Creation, Bounty, and Publishing
        // ================================

        Question draftQuestion = user1.create(
            "What is the difference between interface and abstract class in Java?",
            "I am confused about when to use interface and when to use abstract class. Can someone explain?",
            tags
        );

        System.out.println("=== Scenario 2: Alice drafted a question: ===");
        System.out.println("- Title: What is the difference between interface and abstract class in Java?");
        System.out.println("- Tags: java, oop
");

        user1.addBounty(100);
        System.out.println("Alice added a 100 reputation point bounty to her question.
");

        user1.publishQuestion();
        System.out.println("Alice published the question to the platform.
");

        // ================================
        // 💬 Scenario 3: Answering and Accepting the Answer
        // ================================

        Answer answer = new Answer();
        answer.setId(1);
        answer.setContent("Use interface when you want to define a contract. Use abstract class when you want base functionality.");
        answer.setPostedBy(user2);

        user2.addAnswer(draftQuestion, answer);
        System.out.println("=== Scenario 3: Bob submitted an answer to Alice's question. ===");
        System.out.println("- Answer ID: 1");
        System.out.println("- Content: Use interface when you want to define a contract. Use abstract class when you want base functionality.
");

        user1.acceptAnswer(answer);
        System.out.println("Alice accepted Bob's answer as the correct one.
");

        System.out.println("=== End of Simulation ===");
    }
}
\end{lstlisting}

\subsection{Wrapping up}\label{d6DcwCJ0jKUAiX5iUpOvw-2}
\\
In this chapter, we've explored Stack Overflow's complete design, including how it can be visualized using various UML diagrams and designed using object-oriented principles and design patterns.

% End of section content